On Strata 0.1.27 with Qwen3.8-Flash-Next IQ3\_S, aggregate completion throughput across independent requests increases from \textbf{71.5 tok/s} on one RTX 5070 Ti to \textbf{137.6 tok/s} on two and \textbf{191.8 tok/s} on three, corresponding to point speedups of \textbf{1.93$\times$} and \textbf{2.68$\times$} and parallel efficiencies of \textbf{96.3\%} and \textbf{89.5\%}. Sharing the host expert arena reduces a two-engine PSS snapshot from \textbf{95.3 GiB} to \textbf{52.1 GiB}, a measured reduction of \textbf{43.2 GiB (45.4\%)}; the same 46.8-GiB \texttt{rw-s} mapping is physically shared by both processes. In an interleaved same-generation heterogeneous experiment, an RTX 5070 Ti averages 73.01 tok/s alone and 72.36 tok/s while an RTX 5060 Ti concurrently runs at 59.14 tok/s lane-local TG. The unadjusted fast-lane difference is $-0.65$ tok/s ($-0.9\%$). Because per-run TG is strongly correlated with speculative acceptance ($r=0.93$ solo, $r=0.95$ concurrent), an exploratory acceptance-adjusted OLS estimate is $-1.00$ tok/s ($-1.4\%$, 95\% CI $[-1.71,-0.30]$, $p=0.009$). We interpret this as evidence of a small shared-resource cost, while treating the adjusted estimate as sensitivity analysis rather than a causal decomposition.

The measured operating point is therefore memory-feasible independent serving over shared pinned host state. The evidence supports strong three-lane scaling and physical host-memory sharing on one consumer workstation; broader heterogeneity and direct comparisons with coupled execution remain future work.
\end{abstract}

\section{Introduction}
Large sparse MoE models are attractive on consumer hardware because only a subset of experts is active for each token. Their total expert population can nevertheless exceed the VRAM of a 16~GB GPU. Offload runtimes therefore place some or most expert state in host memory and keep only a working set on the accelerator, turning inference into a hierarchy involving GPU VRAM, host DRAM, PCIe, CPU execution, and sometimes storage~\cite{flexgen,fiddler,ktransformers,moelightning}.

Multiple consumer GPUs create a second constraint: the devices and links are not necessarily symmetric. Coupled multi-GPU strategies may place multiple devices on one request's progress path, and prior heterogeneous-serving systems explicitly model unequal GPU and network capacities~\cite{hexgen,helix,melange}. DWDP likewise identifies layer-wise synchronization and load imbalance as a source of sensitivity in distributed MoE inference~\cite{dwdp}. This paper does \emph{not} measure a coupled baseline on the reference workstation, so we use these systems only as motivation rather than as a performance comparison.

